\documentclass[12pt,a4paper]{article}
\usepackage[utf8]{inputenc}
\usepackage[spanish]{babel}
\usepackage{amsmath}
\usepackage{amssymb}
\usepackage{geometry}
\usepackage{fancyhdr}
\usepackage{titlesec}
\usepackage{xcolor}
\usepackage[hidelinks]{hyperref}
\usepackage{float}
\usepackage{tikz}
\usetikzlibrary{shapes,shapes.geometric,arrows,positioning,babel,backgrounds,calc,decorations.pathmorphing}

% Colores inspirados en pizarrón escolar
\definecolor{chalkboard}{HTML}{1A3020} % Verde oscuro de pizarrón escolar

\tikzset{
    chalkboard style/.style={
        background rectangle/.style={fill=chalkboard},
        show background rectangle,
        inner sep=12pt,
        rounded corners=5pt,
        every label/.style={text=white},
        text=white
    },
    chalkaxis/.style={
        draw=white!85!gray,
        thick,
        ->,
        decorate,
        decoration={random steps, segment length=2mm, amplitude=0.2pt}
    },
    chalkline/.style={
        draw=white!95!gray,
        thick,
        decorate,
        decoration={random steps, segment length=2mm, amplitude=0.15pt}
    },
    chalkdashed/.style={
        draw=white!70!gray,
        dashed,
        thick,
        decorate,
        decoration={random steps, segment length=2.5mm, amplitude=0.2pt}
    },
    chalkvector/.style={
        draw=cyan!90!white,
        ultra thick,
        ->,
        decorate,
        decoration={random steps, segment length=1.5mm, amplitude=0.2pt}
    },
    chalkchord/.style={
        draw=red!80!white,
        ultra thick,
        ->,
        decorate,
        decoration={random steps, segment length=1.5mm, amplitude=0.2pt}
    },
    chalkangle/.style={
        draw=yellow!90!orange,
        thick,
        <->,
        decorate,
        decoration={random steps, segment length=1.5mm, amplitude=0.1pt}
    },
    chalkcircle/.style={
        draw=white!45!gray,
        thin,
        decorate,
        decoration={random steps, segment length=3mm, amplitude=0.25pt}
    }
}

% Configuración de márgenes
\geometry{top=2.5cm, bottom=2.5cm, left=2.5cm, right=2.5cm}

% Encabezado y pie de página
\pagestyle{fancy}
\fancyhf{}
\lhead{\small Especificación Técnica: Mapeo Cartográfico}
\rhead{\small Next.js + SVG}
\rfoot{\thepage}
\lfoot{\small Fundamento Matemático Isótropo}

% Estilo de títulos (Super formal en Blanco y Negro y Negritas)
\titleformat{\section}
  {\normalfont\Large\bfseries}
  {\thesection}{1em}{}
\titleformat{\subsection}
  {\normalfont\large\bfseries}
  {\thesubsection}{1em}{}

% Título del Documento
\title{
    \vspace{2cm}
    \textbf{Especificación Técnica Completa:\\Sistema de Mapeo Cartográfico Local e Interactivo (Next.js + SVG)}\\
    \vspace{1cm}
}
\author{
    \textbf{Autor:} Edgarsc0\\
    \small{Repositorio: \href{https://github.com/Edgarsc0/proyecto_computer_graphics}{Edgarsc0/proyecto\_computer\_graphics}}
}
\date{\today}

\begin{document}

\maketitle
\thispagestyle{empty}
\newpage

Este documento contiene la especificación formal del desarrollo del proyecto, desde el planteamiento del problema hasta su fundamentación matemática e implementación. El objetivo del sistema es proyectar coordenadas geográficas reales (WGS84) en un lienzo bidimensional SVG de forma isótropa y conforme localmente, aplicándolo al campus de UPIITA-IPN.

\section{Introducción}

\subsection{Contexto general}
La representación digital y visualización de espacios geográficos en pantallas bidimensionales es un área fundamental de la geomática y la graficación por computadora. Tradicionalmente, los sistemas cartográficos dependen de estándares complejos como la proyección Web Mercator, los cuales introducen distorsiones significativas a nivel global. Sin embargo, para entornos locales y acotados, tales como un campus universitario, el desarrollo de un motor cartográfico simplificado basado en gráficos vectoriales redimensionables (SVG) provee una solución óptima en rendimiento, flexibilidad y control interactivo sin dependencia de APIs externas.

\subsection{Importancia de la graficación en el problema}
El desafío de renderizar un mapa interactivo radica en la traducción fiel de coordenadas geodésicas tridimensionales (latitud y longitud sobre un elipsoide o esfera terrestre) a coordenadas cartesianas bidimensionales en píxeles del Viewport de la pantalla. Si esta transformación no se calcula correctamente, las proporciones físicas de los edificios, calles y áreas verdes se verán gravemente distorsionadas (achatadas o alargadas). La graficación por computadora aporta los conceptos fundamentales de transformaciones afines, mapeos conformes e isotropía de escala (relación de aspecto métrica uniforme en ambos ejes) para garantizar que la geometría visualizada conserve su forma geométrica real y su orientación local de manera exacta.

\subsection{Descripción breve del proyecto}
Este proyecto consiste en un Sistema de Mapeo Cartográfico Local e Interactivo desarrollado bajo una arquitectura web moderna con Next.js (React) para la visualización interactiva y Django REST Framework para el almacenamiento y gestión de datos geoespaciales. El núcleo del sistema es un motor de transformación geométrica isótropa y conforme que proyecta la geografía del campus escolar en elementos SVG nativos del navegador. La interfaz de usuario permite navegación por paneo, zoom interactivo mediante gestos táctiles (pinch-to-zoom) y edición interactiva de la geometría de los edificios arrastrando vértices en tiempo real, integrando geoperimetraje (geofencing) para el seguimiento de rutas GPS.

\section{Planteamiento del problema}

\subsection{Descripción detallada del problema}
El campus universitario cuenta con diversas infraestructuras cuyas coordenadas espaciales se registran en el formato geodésico estándar WGS84. Para desarrollar aplicaciones de navegación local o de gestión de planta física, se requiere un visor gráfico que muestre este entorno. El problema consiste en diseñar e implementar las ecuaciones de proyección cartográfica local que mapeen estas coordenadas geográficas a un lienzo digital SVG de forma isótropa (donde un metro en el eje horizontal represente exactamente la misma cantidad de píxeles que en el eje vertical) sin recurrir a bibliotecas pesadas de terceros (como Google Maps o Leaflet) y permitiendo la interacción directa (edición de vértices, traslación, escala) sobre los elementos gráficos nativos.

\subsection{Situación actual}
La visualización cartográfica en la web usualmente se soluciona embebiendo mapas propietarios con limitaciones de licencia, o mediante librerías SIG (Sistemas de Información Geográfica) genéricas que consumen gran cantidad de recursos de procesamiento en dispositivos móviles. Además, las aproximaciones lineales simples para la conversión de grados (lat/lon) a píxeles ignoran la curvatura de la Tierra y la reducción del radio del paralelo local con la latitud, lo cual provoca deformaciones severas en la relación de aspecto visual de los edificios.

\subsection{Justificación}
La creación de un motor matemático de mapeo isótropo propio elimina los costos de licenciamiento, optimiza los tiempos de renderizado al interactuar directamente con el árbol DOM de SVG y permite el funcionamiento sin conexión a internet. La justificación técnica radica en resolver la proyección a nivel de álgebra lineal y geometría diferencial para lograr el máximo rendimiento gráfico posible en dispositivos de bajos recursos.

\subsection{Alcances y limitaciones}
\begin{itemize}
    \item \textbf{Alcances:} Mapeo completo y preciso de las poligonales de edificios y senderos del campus; interacción táctil fluida (paneo y zoom mediante gestos); actualización y persistencia en base de datos PostgreSQL de la edición de vértices; y simulación o rastreo GPS en tiempo real con alertas de salida de límites (geofencing).
    \item \textbf{Limitaciones:} La aproximación esférica local y las constantes de transformación son válidas únicamente dentro del polígono delimitador (Bounding Box) del campus (aproximadamente $1\text{ km}^2$), perdiendo validez geométrica a escalas mayores (distancias intermunicipales o globales).
\end{itemize}

\section{Objetivos}

\subsection{Objetivo general}
Desarrollar un sistema interactivo y responsivo de visualización cartográfica local para el campus UPIITA-IPN mediante la implementación de un motor de transformación geodésica isótropa y conforme en Next.js y SVG, conectado a un backend en Django y PostgreSQL.

\subsection{Objetivos específicos}
\begin{itemize}
    \item Deducir las ecuaciones analíticas de transformación entre coordenadas WGS84 y coordenadas de pantalla SVG, calculando la altura exacta del ViewBox para garantizar isotropía.
    \item Implementar la transformación inversa para permitir la recuperación de coordenadas geodésicas a partir de interacciones táctiles en píxeles.
    \item Desarrollar el algoritmo de desplazamiento geodésico en $\mathbb{R}^3$ para proyectar trayectorias y desplazamientos GPS precisos sobre la superficie terrestre.
    \item Construir una interfaz gráfica de alto impacto visual con soporte nativo para gestos móviles multitáctiles (pan, zoom y arrastre de vértices).
\end{itemize}


\section{Fundamento Matemático y Deducción de la Transformación}

\subsection{El Vector de Posición Geocéntrico ($\mathbb{R}^3$)}

La superficie de la Tierra se parametriza en un sistema cartesiano tridimensional ECEF (Earth-Centered, Earth-Fixed) mediante el siguiente vector de posición, donde $r_t$ representa el radio medio terrestre ($\approx 6,371,000\text{ m}$), $\alpha$ la latitud y $\beta$ la longitud:
\begin{equation}
    \mathbf{V}(\alpha,\beta) = r_t \cos(\alpha)\cos(\beta)\mathbf{\hat{i}} + r_t \cos(\alpha)\sin(\beta)\mathbf{\hat{j}} + r_t \sin(\alpha)\mathbf{\hat{k}}
\end{equation}

\begin{figure}[H]
\centering
\begin{tikzpicture}[scale=2.5, >=stealth, chalkboard style,
    axis/.style={chalkaxis},
    vector/.style={chalkvector},
    proj/.style={chalkdashed},
    angle/.style={chalkangle}
]
    % Draw sphere grid (globe representation)
    \draw[chalkcircle] (0,0) circle (1);
    \draw[chalkcircle, dashed] (0,0) ellipse (1 and 0.4);
    
    % Draw axes
    \draw[axis] (0,0,0) -- (1.5,0,0) node[right, color=white] {$x$ (Greenwich, $\beta=0$)};
    \draw[axis] (0,0,0) -- (0,1.5,0) node[above, color=white] {$z$ (Polo Norte)};
    \draw[axis] (0,0,0) -- (0,0,1.8) node[left, color=white] {$y$};
    
    % Draw vector v
    \coordinate (O) at (0,0);
    \coordinate (V) at (0.58,0.57,0.58);
    \coordinate (ProjP) at (0.58,0,0.58);
    
    \draw[vector] (O) -- (V) node[above right, color=white] {$\vec{v} = [v_x, v_y, v_z]^T$};
    \draw[proj] (V) -- (ProjP);
    \draw[proj] (O) -- (ProjP) node[midway, below, color=white] {$\vec{\xi}$};
    \draw[proj] (ProjP) -- (0.58,0,0);
    \draw[proj] (ProjP) -- (0,0,0.58);
    
    % Draw Latitud alpha angle
    \draw[draw=orange!80!white, thick, decorate, decoration={random steps, segment length=1.5mm, amplitude=0.15pt}] (0.2,0,0.2) arc (-45:20:0.28);
    \node[orange!85!white] at (0.32,0.12,0.1) {$\alpha$};
    
    % Draw Longitud beta angle
    \draw[draw=yellow!80!orange, thick, decorate, decoration={random steps, segment length=1.5mm, amplitude=0.15pt}] (0.3,0,0) arc (0:-45:0.3);
    \node[yellow!90!orange] at (0.32,-0.12,0.05) {$\beta$};
\end{tikzpicture}
\caption{Representación Geométrica en $\mathbb{R}^3$ del Vector de Posición Geocéntrico (Estilo Pizarrón).}
\label{fig:geocentrico}
\end{figure}

A partir de este vector, se deduce que las distancias máximas reales del terreno en metros para un área acotada son:
\begin{equation}
    N_{max} = r_t (\alpha_{max} - \alpha_{min})
\end{equation}
para el eje Norte-Sur, y:
\begin{equation}
    E_{max} = r_t \cos(\alpha_0)(\beta_{max} - \beta_{min})
\end{equation}
para el eje Este-Oeste (donde $\alpha$, $\beta$ están expresados en radianes).

\subsection{Dimensión del Viewport ($H$) bajo Condición de Isotropía}

El principio de isotropía exige que la escala del mapa sea idéntica en ambas direcciones de la pantalla para evitar que los objetos se deformen (se estiren o se aplasten). Definimos las escalas en píxeles por metro como:
\begin{equation}
    k_x = \frac{W}{E_{max}}, \quad k_y = \frac{H}{N_{max}}
\end{equation}

Para evitar distorsiones, imponemos la igualdad de escalas $k_x = k_y$:
\begin{equation}
    \frac{W}{E_{max}} = \frac{H}{N_{max}} \implies H = W \cdot \frac{N_{max}}{E_{max}}
\end{equation}

Sustituyendo las distancias métricas del terreno:
\begin{equation}
    H = W \cdot \frac{r_t \cdot (\alpha_{max} - \alpha_{min})}{r_t \cdot \cos(\alpha_0) \cdot (\beta_{max} - \beta_{min})}
\end{equation}

Al simplificar, el radio de la Tierra ($r_t$) se cancela algebraicamente de forma exacta, obteniendo la relación estricta para la altura del viewBox:
\begin{equation}
    H = W \cdot \frac{\alpha_{max} - \alpha_{min}}{(\beta_{max} - \beta_{min}) \cos(\alpha_0)}
\end{equation}

Donde $\alpha_0$ es la latitud central del campus:
\begin{equation}
    \alpha_0 = \frac{\alpha_{max} + \alpha_{min}}{2}
\end{equation}

\subsection{Demostración de las Ecuaciones de la Transformación $T(\beta,\alpha) \to (x,y)$}

Al mapear las distancias a las proporciones de píxeles del lienzo SVG ($W,H$), el radio de la Tierra se anula en ambas componentes:

\begin{enumerate}
    \item \textbf{Componente Horizontal ($x$):} Multiplica la razón de escala del lienzo ($k_x$) por la distancia lineal recorrida desde el límite Oeste ($\beta_{min}$):
    \begin{align}
        x(\beta) &= k_x \cdot \text{Distancia}_E \nonumber \\
        x(\beta) &= \frac{W}{r_t \cos(\alpha_0)(\beta_{max} - \beta_{min})} \cdot \left[ r_t \cos(\alpha_0)(\beta - \beta_{min}) \right] \nonumber \\
        x(\beta) &= W \cdot \frac{\beta - \beta_{min}}{\beta_{max} - \beta_{min}}
    \end{align}

    \item \textbf{Componente Vertical ($y$):} Dado que el eje Y del SVG crece hacia abajo en las pantallas, calculamos la distancia descendente desde el límite Norte ($\alpha_{max}$):
    \begin{align}
        y(\alpha) &= k_y \cdot \text{Distancia}_N \nonumber \\
        y(\alpha) &= \frac{H}{r_t (\alpha_{max} - \alpha_{min})} \cdot \left[ r_t (\alpha_{max} - \alpha) \right] \nonumber \\
        y(\alpha) &= H \cdot \frac{\alpha_{max} - \alpha}{\alpha_{max} - \alpha_{min}}
    \end{align}
\end{enumerate}

\section{Transformación Inversa: $\vec{v} \to (\beta, \alpha)$}

Dado un vector de posición cartesiano tridimensional conocido $\vec{v} = [v_x, v_y, v_z]^T$, es posible recuperar analíticamente las coordenadas geodésicas angulares originales de latitud ($\alpha$) y longitud ($\beta$) mediante las siguientes deducciones geométricas:

\subsection{Cálculo de la Latitud ($\alpha$)}
Despejando directamente a partir de la componente vertical $v_z = r_t \sin \alpha$:
\begin{equation}
    \sin \alpha = \frac{v_z}{r_t} \implies \alpha = \arcsin \left( \frac{v_z}{r_t} \right)
\end{equation}

\subsection{Cálculo de la Longitud ($\beta$)}
Para determinar la longitud, se proyecta el vector de posición $\vec{v}$ sobre el plano horizontal ecuatorial $xy$, obteniendo el vector proyectado $\vec{\xi} = [v_x, v_y, 0]$ cuya magnitud $\xi$ es:
\begin{equation}
    \xi = |\vec{\xi}| = \sqrt{v_x^2 + v_y^2}
\end{equation}
A partir del producto punto entre la proyección $\vec{\xi}$ y el vector unitario director del eje $X$ positivo ($\mathbf{\hat{i}}$), se obtiene la magnitud del ángulo de longitud:
\begin{equation}
    \cos|\beta| = \frac{v_x}{\sqrt{v_x^2 + v_y^2}} \implies |\beta| = \arccos \left( \frac{v_x}{\sqrt{v_x^2 + v_y^2}} \right)
\end{equation}
Dado que el coseno es una función par ($\cos(-x) = \cos x$), se debe resolver la ambigüedad de signo analizando la componente $v_y$ (que define la dirección de la longitud en los hemisferios Este u Oeste):
\begin{equation}
    \beta = \begin{cases} -\arccos \left( \frac{v_x}{\sqrt{v_x^2 + v_y^2}} \right) & \text{si } v_y < 0 \\ \arccos \left( \frac{v_x}{\sqrt{v_x^2 + v_y^2}} \right) & \text{si } v_y \ge 0 \end{cases}
\end{equation}

Esta definición condicional resuelve analíticamente el cuadrante del plano horizontal en el que se ubica el vector proyectado $\vec{\xi}$, mapeando correctamente la longitud $\beta \in (-\pi, \pi]$ de acuerdo con la dirección angular de los vectores en los cuatro cuadrantes (donde $\beta < 0$ si va en sentido de las manecillas del reloj y $\beta \ge 0$ en sentido contrario), tal como se detalla en el siguiente diagrama:

\begin{figure}[H]
\centering
\begin{tikzpicture}[scale=2.0, >=stealth, chalkboard style,
    axis/.style={chalkaxis},
    vector/.style={chalkvector},
    circle_style/.style={chalkcircle},
    angle/.style={chalkangle}
]
    % Draw grid / axes
    \draw[axis] (-2,0) -- (2,0) node[right, color=white] {$x$};
    \draw[axis] (0,-2) -- (0,2) node[above, color=white] {$y$};
    
    % Draw projection circle
    \draw[circle_style] (0,0) circle (1.5);
    \node[white, above right] at (0.9,0.9) {\small Plano $xy$};
    
    % Define vectors in 4 quadrants
    \coordinate (O) at (0,0);
    \coordinate (Q1) at (1.06, 1.06);   % 45 deg
    \coordinate (Q2) at (-1.06, 1.06);  % 135 deg
    \coordinate (Q3) at (-1.06, -1.06); % -135 deg
    \coordinate (Q4) at (1.06, -1.06);  % -45 deg
    
    % Draw vectors
    \draw[vector, draw=cyan!80!white] (O) -- (Q1) node[above right, color=white] {$\vec{\xi}_1$};
    \draw[vector, draw=yellow!80!white] (O) -- (Q2) node[above left, color=white] {$\vec{\xi}_2$};
    \draw[vector, draw=orange!80!white] (O) -- (Q3) node[below left, color=white] {$\vec{\xi}_3$};
    \draw[vector, draw=pink!80!white] (O) -- (Q4) node[below right, color=white] {$\vec{\xi}_4$};
    
    % Draw angles
    % Q1: beta_1 > 0
    \draw[draw=cyan!90!white, thick, ->, decorate, decoration={random steps, segment length=1.5mm, amplitude=0.1pt}] (0.4,0) arc (0:45:0.4);
    \node[cyan!90!white] at (0.6,0.25) {\small $\beta_1 > 0$};
    
    % Q2: beta_2 > 0
    \draw[draw=yellow!90!white, thick, ->, decorate, decoration={random steps, segment length=1.5mm, amplitude=0.1pt}] (0.5,0) arc (0:135:0.5);
    \node[yellow!90!white] at (-0.5,0.7) {\small $\beta_2 > 0$};
    
    % Q3: beta_3 < 0
    \draw[draw=orange!90!white, thick, ->, decorate, decoration={random steps, segment length=1.5mm, amplitude=0.1pt}] (0.6,0) arc (0:-135:0.6);
    \node[orange!90!white] at (-0.5,-0.8) {\small $\beta_3 < 0$};
    
    % Q4: beta_4 < 0
    \draw[draw=pink!90!white, thick, ->, decorate, decoration={random steps, segment length=1.5mm, amplitude=0.1pt}] (0.7,0) arc (0:-45:0.7);
    \node[pink!90!white] at (0.8,-0.4) {\small $\beta_4 < 0$};
\end{tikzpicture}
\caption{Análisis de Signos de la Longitud $\beta$ en los Cuatro Cuadrantes del Plano Ecuatorial (Estilo Pizarrón).}
\label{fig:cuadrantes}
\end{figure}

\section{Algoritmo de Desplazamiento Geodésico Sobre la Esfera}

Dadas dos ubicaciones geodésicas en la Tierra representadas por sus vectores de posición cartesianos conocidos $\vec{v_1}$ y $\vec{v_2}$ (ambos de magnitud $r_t$), se desea calcular el vector de posición $\vec{v_3}$ de una tercera ubicación situada a una distancia real de arco $\phi$ (en metros) sobre la superficie, partiendo desde $\vec{v_1}$ en la dirección de la cuerda recta $\vec{\xi} = \vec{v_2} - \vec{v_1}$ en $\mathbb{R}^3$.

Para resolver esto, modelamos un triángulo plano $OAB$ en el espacio tridimensional (donde $O$ es el origen del planeta, $A$ es el punto inicial de $\vec{v_1}$, y $B$ es el punto auxiliar de desplazamiento lineal a lo largo de la secante):

\begin{figure}[H]
\centering
\begin{tikzpicture}[scale=1.8, >=stealth, chalkboard style,
    origin/.style={circle, draw=white, fill=white, inner sep=1.5pt},
    node/.style={circle, draw=white, fill=cyan!80!white, inner sep=1.5pt},
    chord/.style={chalkchord},
    radius/.style={chalkdashed},
    angle/.style={chalkangle}
]
    % Origin
    \node[origin, label={[text=white]below:$O$}] (O) at (0,0) {};
    
    % Point A (v1)
    \coordinate (A) at (0,4);
    \node[node, label={[text=white]left:{$A (\vec{v_1})$}}] at (A) {};
    \draw[chalkaxis] (O) -- (A) node[midway, left, color=white] {$r_t$};
    
    % Point C (v2)
    \coordinate (C) at (2.5,3.12);
    \node[node, label={[text=white]right:{$C (\vec{v_2})$}}] at (C) {};
    \draw[chalkaxis] (O) -- (C) node[midway, right, color=white] {$r_t$};
    
    % Chord xi
    \draw[chord] (A) -- (C) node[midway, above, color=white] {$\vec{\xi} = \vec{v_2} - \vec{v_1}$};
    
    % Point B (intermediate w)
    \coordinate (B) at (1.2,3.58);
    \node[circle, draw=white, fill=pink!80!white, inner sep=1.2pt, label={[text=white]above right:{$B (\vec{w})$}}] at (B) {};
    \draw[chalkdashed] (O) -- (B);
    
    % Intersection point v3
    \coordinate (V3) at (0.79, 3.92); % projected on sphere
    \node[circle, draw=white, fill=green!80!white, inner sep=1.5pt, label={[text=white]above:{$v_3$}}] at (V3) {};
    
    % Draw earth surface arc
    \draw[draw=cyan!80!white, thick, decorate, decoration={random steps, segment length=2mm, amplitude=0.15pt}] (0,4) arc (90:51.3:4);
    \node[cyan!80!white] at (1.5, 3.8) {Arco $\phi$};
    
    % Draw angle theta at O
    \draw[draw=yellow!80!orange, thick, decorate, decoration={random steps, segment length=1.5mm, amplitude=0.1pt}] (0,0.8) arc (90:51.3:0.8);
    \node[yellow!80!orange] at (0.3,1.0) {$\theta$};
    
    % Draw angle gamma2 and gamma1 at A
    \draw[draw=orange!80!white, thick, decorate, decoration={random steps, segment length=1mm, amplitude=0.1pt}] (0,3.5) arc (270:345:0.5);
    \node[orange!80!white] at (0.2,3.3) {$\gamma_2$};
    
    \draw[draw=pink!80!white, thick, decorate, decoration={random steps, segment length=1mm, amplitude=0.1pt}] (0,3.4) arc (270:90:0.6);
    \node[pink!80!white] at (-0.3,3.4) {$\gamma_1$};
    
    % Draw angle sigma at B
    \draw[draw=yellow!80!orange, thick, decorate, decoration={random steps, segment length=1mm, amplitude=0.1pt}] (1.2,3.58) -- ++(-0.2,-0.6);
    \node[yellow!80!orange] at (0.9,3.3) {$\sigma$};
\end{tikzpicture}
\caption{Diagrama del Triángulo Auxiliar $OAB$ para la Deducción del Desplazamiento Geodésico (Estilo Pizarrón).}
\label{fig:desplazamiento}
\end{figure}

\subsection{Cálculo de la Dirección y el Ángulo Central ($\theta$)}
Definimos el vector unitario director de la cuerda recta en el espacio:
\begin{equation}
    \mathbf{\hat{\xi}} = \frac{\vec{v_2} - \vec{v_1}}{|\vec{v_2} - \vec{v_1}|}
\end{equation}
El desplazamiento métrico de arco $\phi$ equivale a un ángulo central $\theta$ (en grados) de:
\begin{equation}
    \theta = \frac{180 \phi}{\pi r_t}
\end{equation}

\subsection{Cálculo de los Ángulos Internos del Triángulo $OAB$}
\begin{enumerate}
    \item \textbf{Ángulo entre vectores $\vec{v_1}$ y $\vec{\xi}$ ($\gamma_2$):} Se calcula mediante la definición geométrica del producto punto:
    \begin{equation}
        \vec{v_1} \cdot \vec{\xi} = r_t |\vec{\xi}| \cos \gamma_2 \implies \gamma_2 = \arccos \left( \frac{\vec{v_1} \cdot \vec{\xi}}{r_t |\vec{\xi}|} \right)
    \end{equation}
    \item \textbf{Ángulo interior suplementario ($\gamma_1$):} Dado que $\gamma_1$ y $\gamma_2$ son suplementarios:
    \begin{equation}
        \gamma_1 = 180^\circ - \gamma_2 = 180^\circ - \arccos \left( \frac{\vec{v_1} \cdot \vec{\xi}}{r_t |\vec{\xi}|} \right)
    \end{equation}
    \item \textbf{Ángulo opuesto en el vértice $B$ ($\sigma$):} Por suma de ángulos internos de un triángulo plano ($180^\circ$):
    \begin{align}
        \sigma &= 180^\circ - \gamma_1 - \theta \nonumber \\
        \sigma &= \arccos \left( \frac{\vec{v_1} \cdot \vec{\xi}}{r_t |\vec{\xi}|} \right) - \theta
    \end{align}
\end{enumerate}

\subsection{Escalamiento Lineal de la Cuerda ($k_0$) y Proyección Final}
Aplicando la Ley de Senos sobre el triángulo plano $OAB$:
\begin{equation}
    \frac{r_t}{\sin \sigma} = \frac{k_0}{\sin \theta} \implies k_0 = r_t \frac{\sin \theta}{\sin \sigma}
\end{equation}
Sustituyendo los valores angulares deducidos, obtenemos el factor de escala lineal $k_0$:
\begin{equation}
    k_0 = r_t \frac{\sin \left( \frac{180 \phi}{\pi r_t} \right)}{\sin \left( \arccos \left( \frac{\vec{v_1} \cdot \vec{\xi}}{r_t |\vec{\xi}|} \right) - \frac{180 \phi}{\pi r_t} \right)}
\end{equation}
Calculamos el vector secante intermedio $\vec{w}$ desplazado a lo largo de la cuerda:
\begin{equation}
    \vec{w} = \vec{v_1} + k_0 \mathbf{\hat{\xi}}
\end{equation}
Finalmente, normalizamos y proyectamos el vector de vuelta a la superficie esférica terrestre (alargándolo al radio $r_t$):
\begin{equation}
    \vec{v_3} = r_t \mathbf{\hat{w}} = r_t \frac{\vec{w}}{|\vec{w}|}
\end{equation}
Como se detalla en los diagramas del pizarrón escolar original, esta proyección radial de vuelta a la superficie esférica se puede expresar de forma alternativa sumando al vector secante $\vec{w}$ el vector del segmento radial restante $(r_t - |\vec{w}|)$ en su propia dirección unitaria $\mathbf{\hat{w}}$:
\begin{equation}
    \vec{v_3} = \vec{w} + (r_t - |\vec{w}|) \mathbf{\hat{w}}
\end{equation}
donde $\mathbf{\hat{w}} = \frac{\vec{w}}{|\vec{w}|}$. Esta formulación hace más intuitiva la noción física de alargar radialmente el punto secante $B(\vec{w})$ hasta intersectar la superficie esférica de la Tierra en $v_3$.

\section{Universo Geográfico del Campus (Bounding Box)}

Los límites absolutos extraídos del polígono envolvente de la escuela definen las constantes de la transformación:
\begin{itemize}
    \item \textbf{Latitud Máxima ($\alpha_{max}$):} $19.506085642577684^\circ$
    \item \textbf{Latitud Mínima ($\alpha_{min}$):} $19.501672871991357^\circ$
    \item \textbf{Longitud Máxima ($\beta_{max}$):} $-99.14538442387669^\circ$
    \item \textbf{Longitud Mínima ($\beta_{min}$):} $-99.14882990272824^\circ$
\end{itemize}

\subsection{Parámetros Calculados del Lienzo}
\begin{itemize}
    \item $\Delta\alpha = 0.004412770586327^\circ$
    \item $\Delta\beta = 0.003445478851546^\circ$
    \item $\alpha_0 = 19.50387925728452^\circ \implies \cos(\alpha_{0,\text{rad}}) \approx 0.942620925$
\end{itemize}

Fijando un ancho de $W=1000$, la altura calculada bajo isotropía es $H=1358.71$. El contenedor debe inicializarse con el atributo:
\begin{center}
    \texttt{viewBox="0 0 1000 1358.71"}
\end{center}

\end{document}
